
We investigated whether neural networks can learn permutation invariance from data diversity or whether such structure must be built into the architecture. In weight-space learning for accuracy prediction, MLPs achieve near-perfect predictive performance ($R^{2}=0.986)$ when trained on 40,000 models but fail to develop permutation invariance (score 0.63 versus NFN's 1.0).

This demonstrates that task performance and representation quality are separable---a model can learn the right answers through the wrong mechanism. The diagnostic tools developed (permutation invariance testing) and the architectural solution (NFN) ensure that weight-space models not only predict correctly but represent the input structure correctly.

Permutation equivariance must be built into the architecture. Data quantity cannot substitute.

